% !PRINTED 239
\weberpage{239}
Dazu ist erforderlich
\weq{11}{\begin{gathered}
\alpha a_2 + \beta c_2 = 0 \\
\gamma a_2 + \delta c_2 = -\partial,
\end{gathered}}
\weq{12}{\begin{gathered}
\beta \partial_2 = a \\
\delta \partial_2 = c.
\end{gathered}}
Es ist also $\partial_2$ bestimmt als der größte gemeinschaftliche
Teiler von $a$ und $c$, und damit zugleich, wegen $a_2\partial_2 = n$, auch $a_2$.

Nach den beiden Gleichungen (12) kennt man jetzt $\beta$, $\delta$
als relative Primzahlen und kann $\alpha$, $\gamma$ aus der diophantischen
Gleichung
\[
\alpha\delta - \beta\gamma = 1
\]
bestimmen. Ist dies geschehen, so folgt aus den beiden
Gleichungen (11)
\weq{13}{\begin{gathered}
a_2 = \partial\beta, \\
c_2 = -\partial\alpha,
\end{gathered}}
wodurch auch $c_2$ bestimmt ist, und es zeigt sich zugleich, daß $\partial$
der größte gemeinschaftliche Teiler von $a_2$, $c_2$ ist. Ersetzt man
$\alpha$, $\gamma$ durch eine andere Lösung $\alpha + h\beta$, $\gamma + h\delta$, so ändert sich
nur $c_2$ um ein Vielfaches von $a_2$.

Durch die Transformation (7) geht also
\[
j\left(\frac{c+\partial\omega}{a}\right) \text{ in } j\left(\frac{c_2+\partial_2\omega}{a_2}\right)
\]
über.

Bilden wir nun eine \gesperrt{symmetrische} Funktion der sämtlichen
Größen (1), etwa für ein unbestimmtes $x$ das Produkt
\[
\Pi\left[x - j\left(\frac{c+\partial\omega}{a}\right)\right],
\]
so ändert sich diese Funktion nicht durch die Transformationen
(6), (7) und ist also nach dem oben erwähnten Satz eine \gesperrt{rationale
Funktion} von $j(\omega)$. Außerdem ist sie eine ganze rationale
Funktion $\nu$ten Grades von $x$ mit dem Anfangsglied $x^{\nu}$, und
wir bezeichnen sie mit $F_n[x, j(\omega)]$. Die Gleichung
\weq{14}{F_n[x, j(\omega)] = 0}
hat die Größen $j\left(\dfrac{c+\partial\omega}{a}\right)$ zu Wurzeln und ist die \gesperrt{Invariantengleichung},
deren wichtigste Eigenschaften wir nun ableiten wollen.

1. Die Invariantengleichung ist \gesperrt{irreducibel}, wenn als Rationalitätsbereich der Inbegriff aller rationalen Funktionen von
%CONT>>
